%*******************************************************
% Abstract in German
%*******************************************************
\begin{otherlanguage}{ngerman}
    \pdfbookmark[0]{Zusammenfassung}{Zusammenfassung}
    \chapter*{Zusammenfassung}

    Abstrakte Argumentationssysteme modellieren den Umgang mit
    widersprüchlichen Informationen durch Argumente und gerichtete Angriffe.
    Eine Semantik legt fest, welche Gesamtbewertungen der Argumente zulässig
    sind; Labellings stellen diese Bewertungen dar, indem sie jedes Argument
    als akzeptiert, abgelehnt oder unentschieden kennzeichnen. Zwei
    Argumentmengen sind bedingt unabhängig, wenn sich ihre jeweiligen
    Bewertungen aus zwei beliebigen zulässigen Labellings, die auf einer
    dritten Argumentmenge übereinstimmen, in einem weiteren zulässigen
    Labelling zusammenführen lassen. Die Bewertung der dritten Menge bildet
    den Kontext und bleibt dabei erhalten.

    Die theoretischen Eigenschaften bedingter Unabhängigkeit in solchen
    Systemen sind bislang nur teilweise untersucht. Bestehende graphbasierte
    Kriterien erfassen zudem nur einen Teil der geltenden Unabhängigkeiten.
    Diese Arbeit untersucht die rechnerische Schwierigkeit des Problems und
    entwickelt ein exaktes Entscheidungsverfahren.

    Für die häufig untersuchten Semantiken \emph{complete}, \emph{stable}
    und \emph{preferred} werden Unabhängigkeitsanfragen einheitlich in
    quantifizierte boolesche Formeln (QBFs) übersetzt und mit dem QBF-Solver
    QuAbS ausgewertet. Ein Korrektheitsbeweis zeigt, dass eine erzeugte Formel
    genau dann wahr ist, wenn die angefragte Unabhängigkeit gilt.

    Die Komplexitätsanalyse zeigt, dass in Systemen ohne gerichtete Zyklen
    gerader Länge jede Unabhängigkeitsanfrage positiv beantwortet wird. In
    Systemen ohne gerichtete Zyklen ungerader Länge bleibt das Problem unter
    allen drei Semantiken dagegen rechnerisch schwierig
    (\(\Pi_2^{\mathrm{P}}\)-vollständig).

    In einer empirischen Evaluation mit \(9\,450\) Anfragen auf \(105\)
    Systemen lassen sich rund drei Viertel der Anfragen entscheiden. Davon
    bestätigen etwa \(98\,\%\) bedingte Unabhängigkeit. Größere Kontextmengen
    sind unter den entschiedenen Anfragen mit weniger beobachteter
    Abhängigkeit verbunden. Zugleich gehen sie tendenziell mit längeren
    Laufzeiten und einem geringeren Anteil entschiedener Anfragen einher.
\end{otherlanguage}
